\clearpage
\section{Whole-graph reconstruction and binding-free operational residuals}
\label{sec:copying-operational-transfer}
Micro-beta and reference creation have multi-edge left-hand graphs. A
primitive match in one filled specimen therefore cannot be transported by
assuming that every hole contributes a single corolla. This section proves
the missing constructive ingredients: finite complements of complete
matches, actual whole-graph rewrite realization, binding-free residual
certificates, and the copying overlap barrier. The original compute rules,
stable arguments and translation are unchanged.

These results serve the operational-template argument in
\href{https://arxiv.org/html/1907.01257v5\#A8.SS3.SSS7}{Appendix H.3.7}.
The complement theorem constructs a real graph context and proves its
reconstruction isomorphism. The residual results check binding freedom of
the shared context itself. Neither a matching predicate nor an unrestricted
residual specimen is treated as a proof of binding-free robustness.
The universal forward/converse robustness theorem for the two operational
families remains open: every original source overlap must still be connected
to the appropriate local diagram, and the grafted residual context must
be proved binding-free. The generic matched transfer below now constructs
the paired complement and its original recursive residual certificate.

\subsection{A complete match determines its finite complement}
\termdef{complete-match-complement}
A \emph{complete match} consists of an incidence-preserving occurrence
$f:A\to G$ together with recursive agreement of every matched edge label.
Its boundary inputs and outputs may be identified by attachment, but no two
inputs, no two outputs, and no internal vertex with another vertex can be
identified. To prove these facts, use source or target coverage to find
incidences at the two candidate vertices. Incidence uniqueness identifies
their edges; injectivity of the edge map and the absence of duplicate ports
in the ordered incidence list then identify the vertices.

An internal vertex of $A$ has both a producer and a consumer. Its image
therefore cannot be an input or output of $G$. This gives an injective map
between the internal carriers. Recursive label agreement supplies equality
of the edge signatures, so the ordered incidence equations recover every
wire type, including at boundary ports. These facts are consequences of the
match, not additional requirements imposed on the machine.

Remove the images of the internal vertices and edges of $A$ from $G$.
The remaining internal carrier is the subtype of vertices outside the
internal image. The remaining edges retain their labels and ordered
incidences. An incidence of a remaining edge cannot meet a removed internal
vertex: otherwise source or target uniqueness would identify that edge
with a matched edge. Thus these lists project to the remaining carrier.
Insert one new hole edge with the ordered input and output images of $A$.
The resulting context is finite because its carriers are subtypes of the
original finite carriers, with one additional edge.

Coverage and uniqueness are checked separately for retained edges and the
new hole. At a surviving vertex, an original incidence is either retained,
or comes from a matched boundary incidence and is replaced by the hole.
Matched internal incidences cannot occur there. Same-polarity injectivity
proves that each hole incidence list has no repetitions. An input-output
alias is permitted and remains a loop in the hole attachment. The original
wire labels prove typing of every retained incidence and every hole port.

\Needspace{11\baselineskip}
\subsection{Reconstruction and actual rewriting}
\termdef{complete-match-realization}
The \emph{realization isomorphism} identifies the filled complement with
$G$. For any injection, the disjoint sum of its complement and its source
is bijective with its target. Apply this fact to the internal carrier and
to the edge carrier. The resulting maps send retained objects to themselves
and inserted objects to their matched images. Their label and ordered
incidence equations follow from the original match. In particular, the
isomorphism sends every inserted port to $f$ of that port; it does not
merely compare the unpointed graphs.

Given an original rewrite $A,v\longrightarrow B,w$, cast its boundary
indices to the proved hole signature, plug it into the constructed
complement, and transport its source through the realization isomorphism.
This produces an original whole-graph rewrite from $G$ at $f(v)$. Its target
is the same complement filled by $B$, and its resumed token is the actual
inserted image of $w$. Every constructor used here already belongs to the
original rewrite relation.

Conversely, induction on an original contextual rewrite extracts an actual
local rule and an occurrence preserving full recursive labels at the exact
source token. Local steps use the identity occurrence; renaming, plugging
and expansion compose it with their proved occurrence maps. Consequently
the label hypothesis of reconstruction is available for every original
rewrite, including beta and copying rules that inspect boxes. Comparing the supplied
successor with the reconstructed successor uses the already proved
whole-machine successor coherence. The resulting pointed recursive
isomorphism preserves the exact remaining successful-run lengths.

The auxiliary cut hole need not be a binding-free observation hole, and its
name is not asserted fresh in a paper replacement tuple. Reconstruction is
an edge-indexed graph statement. The original-hole inventory correspondence
and admissibility of a paired residual context remain separate obligations.

\subsection{Paired reconstruction in the original common context}
\termdef{paired-footprint-frame}
For a complete match that avoids shallow holes, replace only its auxiliary
edge labels by the corresponding original context labels. The carriers,
ordered incidences and boundary are unchanged. This \emph{matched skeleton}
has a full occurrence in the original common context. Restricting the
original recursive filling certificates to its edges reconstructs the
unchanged actual local source on the left, and an actual paired footprint
on the right. The source-token reconstruction equation is proved explicitly.

Cut this skeleton out of the common context. At the selected cut component,
place its expansion by the chosen original filling family; at every
unmatched component, keep that original filling unchanged. The resulting
whole expansion is recursively isomorphic to the original filled host.
The proof partitions the edge carrier and regroups the internal carrier;
its selected and retained insertion equations verify both ordered
incidences. This works for an arbitrary original filling family, so the
two sides share one computed frame even when unrelated holes have many
internal vertices and edges.

Now supply the independently established two local rewrites and their local
paired residual certificate. Expand those actual rewrites in the common
frame. Use witnesses for every original unmatched endpoint pair and put the local residual
at the selected component. This constructs the two whole-graph rewrites and
a genuine recursive whole-graph filling witness. The left token is the
actual matched token; the right token is the paired host token when the
stated attachment equation holds. An independently supplied left successor
is compared to the constructed one by proved coherence. The final residual
context need not inherit binding freedom merely because the cut exists;
that remaining admissibility proof is not hidden in this construction.
The edge-witness extraction may rebuild its common context and hole names.
Original-hole inventory and pointed residual-coordinate correspondence
must therefore be checked separately for this general grafted certificate.

\subsection{Binding freedom excludes shallow copying overlaps}
\termdef{binding-overlap-barrier}
At an original token that is not entering a hole, an inspected primitive
must lie in a retained context contribution. A purported filled-hole
contribution would place that token at the hole source, contradicting the
non-entering premise. This argument permits arbitrary multi-edge graphs
at all other holes.

A retained predecessor exposes its actual context target ports. Source
ownership reflects the next expanded adjacency to a context adjacency.
Follow an expanded path starting from a sourced binding edge. If it first
enters a shallow hole, its already-reflected prefix gives a positive
context path from that binding edge to the hole. Binding freedom forbids
this path. Induction therefore retains and reflects the entire path.

For copying, the active contraction is such a binding edge. Each sharing
hop reflects its primitive label and adjacency; induction preserves the
exact number of hops. The resolved payload head is consequently present
in the common context at the same cost. Every outer edge of a normalized
copyable payload is either its primitive head or, for a lambda, its
immediately following box. Both are descendants of the contraction and
cannot belong to a shallow hole. This proof leaves deep occurrences inside
the box available for the residual construction.

More generally, when every edge of a matched footprint avoids shallow
holes, incidence coverage places each matched port at an original context
coordinate. Reflect the ordered match to that context, then insert the
opposite retained edge corollas to obtain the opposite match at the exact
paired token. Unrelated hole fillings can have arbitrary internal vertices.
For a primitive footprint, primitive agreement is already full label
agreement, and the complement theorem constructs a genuine opposite-side
rewrite. For a footprint containing boxes, the opposite local source must
instead be rebuilt with its paired bodies; primitive agreement alone does
not supply that construction.

\subsection{Stable arguments and binding-free residual certificates}
\termdef{binding-free-filling-witness}
A \emph{binding-free filling witness} retains the original common context,
its two recursive edge-certificate families and the endpoint isomorphisms,
and additionally proves binding freedom of that common context. Closing
its endpoints and supplying a typed token constructs an actual specimen
in the binding-free context class. For a non-rewrite token the pointed
endpoint equalities give the equality-up-to clause used in robustness.
The required rootedness premises are those of that clause, not assumptions
of termination.

Padding preserves the two original operational pairs: pad their bodies,
arguments and endpoint isomorphisms using the proved compute-padding laws.
Neither side of either pair has exclusively flat stable labels. The left
has its active application or reference head. A beta reduct still contains
the opened body with a term output, impossible for a flat stable graph;
a reference reduct contains the created store. Hence a stable paired
argument has no operational filling and no boxed contribution. All its
certificates are retained corollas, giving an actual recursive isomorphism
between its two endpoints.

Serial and parallel residual assembly retain their actual shared contexts.
For copying a lambda, the outer lambda-and-box context has no shallow
holes; its holes occur only inside the original box body. This permits the
residual contraction tree to precede the retained copy without introducing
a forbidden path to a shallow hole. Both copies retain their certificates. Appending the focus-free, hole-free
distributors preserves binding freedom.

Padding is constructed together with a recursive isomorphism from its
actual shared context to the padded original context. Transporting binding
freedom through this isomorphism justifies opening the padded paired body;
no property of an unconstrained chosen witness is assumed.

For beta opening, the bound body output connects to a closed argument.
A path can move from the body into that argument, but cannot return: the
argument outputs are external outputs of the splice, whereas the body
input is its external input. A path ending at a body hole consequently
stays in the body. A forbidden binding path would contradict the body's
binding freedom; the argument contains no holes. This proves binding
freedom of the actual splice, without a stability or rootedness premise.

To place a residual in a closing host, use an original sourced hole.
The plugged certificate's common context is explicitly isomorphic to the
host plugged by the residual context. Proved sourced-hole closure gives
binding freedom. Expansion peeling identifies both pointed endpoint
states, including internal host resume coordinates, and constructs the
required original binding-free residual specimen.

These constructions prove both forward and converse quantitative robustness
clauses for the framed deep-copying and beta cases. A supplied successor
may use an independent representation: proved successor coherence compares
it to the constructed successor. The matched cost comparison is one with
one, with zero additional transitions after the supplied step. The converse
exchanges the same residual specimen's two projections. The universal
operational-template theorem must still extract these framed data from
every original specimen and handle the remaining overlaps.

\subsection{Concrete nested operational regressions}
The reference regression uses the original constant translation and
allocation graphs. The beta regression uses the unchanged variable body,
including its single-use distributor, and the original stable constant
argument. Tensoring with an actual weakening exposes an unused bound output;
this makes each graph a lambda body without changing its operation.
The copying certificates retain the operational occurrence in both actual
residual branches. A second enclosing lambda checks a further level of
nested boxes. No body-termination premise is used, and copying does not
evaluate the reference operation hidden in a box.

% Generated declaration links; prose states each construction and proof family.
\subsection{Lean declarations and proof organisation}
\Needspace{8\baselineskip}
\paragraph{Port geometry.}
The two same-polarity injectivity lemmas use incidence coverage, edge
injectivity and ordered-list nonduplication. The existence lemma places each
internal image in the host internal carrier; \texttt{internal} chooses that
carrier element and \texttt{ports\_internal} proves its defining equation.
The internal-equality equivalence excludes every other port from that image,
and hence proves internal injectivity. The wire-label theorem uses the
source equation for internal/input ports and the target equation for output
ports, together with the full matched label signatures.

\lean{Copying/MatchEmbeddingGeometry}{25}{33}{MatchEmbedding.ports_injective_of_not_output}
\lean{Copying/MatchEmbeddingGeometry}{36}{44}{MatchEmbedding.ports_injective_of_not_input}
\lean{Copying/MatchEmbeddingGeometry}{47}{56}{MatchEmbedding.exists_internal}
\lean{Copying/MatchEmbeddingGeometry}{59}{60}{MatchEmbedding.internal}
\lean{Copying/MatchEmbeddingGeometry}{62}{64}{MatchEmbedding.ports_internal}
\par\penalty0
\lean{Copying/MatchEmbeddingGeometry}{67}{83}{MatchEmbedding.ports_eq_internal_iff}
\lean{Copying/MatchEmbeddingGeometry}{85}{89}{MatchEmbedding.internal_injective}
\lean{Copying/MatchEmbeddingGeometry}{92}{116}{MatchEmbedding.vertexLabel}

\Needspace{8\baselineskip}
\paragraph{Complement carriers.}
The complement bijection sends an outside element to itself and a source
element to its injected image. Its inverse tests membership in that image;
in the positive case injectivity identifies the chosen preimage, and in the
negative case the point carries its complement proof. The two displayed
summand equations are definitional.

\lean{Graph/ComplementBijection}{9}{34}{Bijection.ofComplement}
\lean{Graph/ComplementBijection}{36}{38}{Bijection.ofComplement_inl}
\lean{Graph/ComplementBijection}{40}{41}{Bijection.ofComplement_inr}

\Needspace{8\baselineskip}
\paragraph{The constructed cut context.}
\texttt{Outside} excludes every matched internal image. The internal,
retained-edge and vertex types are the corresponding finite subtypes and
port type. The embedding forgets subtype proofs; projection adds the proved
outside predicate. Their inverse equations and embedding injectivity follow
by port cases. The source/target outside lemmas use incidence uniqueness;
the retained lists map the original attached lists through projection.
Mapping them back recovers the original lists, giving the membership and
nonduplication lemmas. Projection and embedding preserve wire labels by
port cases.

The input/output outside lemmas use same-polarity injectivity against an
internal port. Their projected ports embed to the original matched ports
and are injective on each boundary separately. The cut source, target and
label functions distinguish a retained edge from the new hole. The two hole
membership lemmas are membership in mapped ordered boundary lists.
Nonduplication follows from the boundary injections; disjointness from a
retained edge follows from original incidence uniqueness. These prove
source and target uniqueness for all four combinations of edge kind.
Coverage uses original coverage, testing whether its edge is in the match:
a matched preimage at a surviving vertex must be the appropriate boundary
port. Typing uses the projection-label equation for retained edges and the
full match's vertex-label theorem for the hole. Finally \texttt{net}
packages precisely these laws and finite carrier enumerations.

\lean{Copying/CompleteMatchCut}{12}{12}{CompleteMatchCut.Outside}
\lean{Copying/CompleteMatchCut}{14}{14}{CompleteMatchCut.Internal}
\lean{Copying/CompleteMatchCut}{15}{15}{CompleteMatchCut.RetainedEdge}
\lean{Copying/CompleteMatchCut}{16}{16}{CompleteMatchCut.Vertex}
\lean{Copying/CompleteMatchCut}{18}{18}{CompleteMatchCut.embed}
\par\penalty0
\lean{Copying/CompleteMatchCut}{20}{24}{CompleteMatchCut.project}
\lean{Copying/CompleteMatchCut}{26}{27}{CompleteMatchCut.embed_project}
\lean{Copying/CompleteMatchCut}{29}{43}{CompleteMatchCut.embed_outside}
\lean{Copying/CompleteMatchCut}{45}{46}{CompleteMatchCut.embed_injective}
\lean{Copying/CompleteMatchCut}{48}{49}{CompleteMatchCut.project_embed}
\par\penalty0
\lean{Copying/CompleteMatchCut}{51}{54}{CompleteMatchCut.outside_source}
\lean{Copying/CompleteMatchCut}{56}{59}{CompleteMatchCut.outside_target}
\lean{Copying/CompleteMatchCut}{61}{62}{CompleteMatchCut.retainedSource}
\lean{Copying/CompleteMatchCut}{64}{65}{CompleteMatchCut.retainedTarget}
\lean{Copying/CompleteMatchCut}{67}{69}{CompleteMatchCut.embed_retainedSource}
\par\penalty0
\lean{Copying/CompleteMatchCut}{71}{73}{CompleteMatchCut.embed_retainedTarget}
\lean{Copying/CompleteMatchCut}{75}{82}{CompleteMatchCut.mem_retainedSource}
\lean{Copying/CompleteMatchCut}{84}{91}{CompleteMatchCut.mem_retainedTarget}
\lean{Copying/CompleteMatchCut}{93}{94}{CompleteMatchCut.nodup_of_map}
\lean{Copying/CompleteMatchCut}{96}{97}{CompleteMatchCut.retainedSource_nodup}
\par\penalty0
\lean{Copying/CompleteMatchCut}{99}{100}{CompleteMatchCut.retainedTarget_nodup}
\lean{Copying/CompleteMatchCut}{102}{104}{CompleteMatchCut.project_label}
\lean{Copying/CompleteMatchCut}{106}{108}{CompleteMatchCut.embed_label}
\lean{Copying/CompleteMatchCut}{110}{114}{CompleteMatchCut.outside_input}
\lean{Copying/CompleteMatchCut}{116}{120}{CompleteMatchCut.outside_output}
\par\penalty0
\lean{Copying/CompleteMatchCut}{122}{122}{CompleteMatchCut.inputPort}
\lean{Copying/CompleteMatchCut}{123}{123}{CompleteMatchCut.outputPort}
\lean{Copying/CompleteMatchCut}{125}{126}{CompleteMatchCut.embed_input}
\lean{Copying/CompleteMatchCut}{127}{128}{CompleteMatchCut.embed_output}
\lean{Copying/CompleteMatchCut}{130}{135}{CompleteMatchCut.inputPort_injective}
\par\penalty0
\lean{Copying/CompleteMatchCut}{137}{142}{CompleteMatchCut.outputPort_injective}
\lean{Copying/CompleteMatchCut}{144}{146}{CompleteMatchCut.source}
\lean{Copying/CompleteMatchCut}{148}{150}{CompleteMatchCut.target}
\lean{Copying/CompleteMatchCut}{152}{154}{CompleteMatchCut.label}
\lean{Copying/CompleteMatchCut}{156}{157}{CompleteMatchCut.mem_hole_source}
\par\penalty0
\lean{Copying/CompleteMatchCut}{159}{160}{CompleteMatchCut.mem_hole_target}
\lean{Copying/CompleteMatchCut}{162}{165}{CompleteMatchCut.source_nodup}
\lean{Copying/CompleteMatchCut}{167}{170}{CompleteMatchCut.target_nodup}
\lean{Copying/CompleteMatchCut}{172}{178}{CompleteMatchCut.source_disjoint}
\lean{Copying/CompleteMatchCut}{180}{186}{CompleteMatchCut.target_disjoint}
\par\penalty0
\lean{Copying/CompleteMatchCut}{188}{198}{CompleteMatchCut.source_unique}
\lean{Copying/CompleteMatchCut}{200}{210}{CompleteMatchCut.target_unique}
\lean{Copying/CompleteMatchCut}{212}{240}{CompleteMatchCut.source_coverage}
\lean{Copying/CompleteMatchCut}{242}{270}{CompleteMatchCut.target_coverage}
\lean{Copying/CompleteMatchCut}{276}{295}{CompleteMatchCut.source_typed}
\par\penalty0
\lean{Copying/CompleteMatchCut}{297}{316}{CompleteMatchCut.target_typed}
\lean{Copying/CompleteMatchCut}{320}{341}{CompleteMatchCut.net}

\Needspace{8\baselineskip}
\paragraph{Reconstruction maps.}
The internal-complement bijection identifies the port-based outside
predicate with the complement of the induced internal injection. Combining
it with the complement bijection gives \texttt{vertices}. The retained-edge
bijection removes the new hole, and the corresponding disjoint-sum
construction gives \texttt{edges}. Their four summand equations specify
exactly where each retained or inserted object goes. The hole input/output
lemmas prove the original ordered interface; \texttt{replacement} performs
only those boundary casts. Host and source/target-port equations reduce
reconstruction on attached vertices to the original match. The insertion
equation handles all three port kinds. The final recursive isomorphism
checks labels and both ordered incidence lists on each edge summand.

\lean{Copying/CompleteMatchRealization}{13}{20}{CompleteMatchCut.internalComplement}
\lean{Copying/CompleteMatchRealization}{24}{26}{CompleteMatchCut.vertices}
\lean{Copying/CompleteMatchRealization}{28}{28}{CompleteMatchCut.vertices_left}
\lean{Copying/CompleteMatchRealization}{29}{30}{CompleteMatchCut.vertices_right}
\lean{Copying/CompleteMatchRealization}{33}{41}{CompleteMatchCut.retainedEdges}
\par\penalty0
\lean{Copying/CompleteMatchRealization}{44}{47}{CompleteMatchCut.edges}
\lean{Copying/CompleteMatchRealization}{49}{50}{CompleteMatchCut.edges_left}
\lean{Copying/CompleteMatchRealization}{51}{51}{CompleteMatchCut.edges_right}
\lean{Copying/CompleteMatchRealization}{55}{56}{CompleteMatchCut.hole_input}
\lean{Copying/CompleteMatchRealization}{58}{59}{CompleteMatchCut.hole_output}
\par\penalty0
\lean{Copying/CompleteMatchRealization}{63}{63}{CompleteMatchCut.replacement}
\lean{Copying/CompleteMatchRealization}{65}{68}{CompleteMatchCut.map_host}
\lean{Copying/CompleteMatchRealization}{70}{73}{CompleteMatchCut.hole_sourcePort}
\lean{Copying/CompleteMatchRealization}{75}{78}{CompleteMatchCut.hole_targetPort}
\lean{Copying/CompleteMatchRealization}{82}{97}{CompleteMatchCut.map_insert}
\par\penalty0
\lean{Copying/CompleteMatchRealization}{101}{151}{CompleteMatchCut.realization}

\Needspace{8\baselineskip}
\paragraph{Actual rewrite realization.}
\texttt{replaceWith} fills the constructed cut with the successor graph;
\texttt{replacePort} is its actual resumed port. The rewrite theorem casts
the proved local interface, uses the original plug constructor, and then
renames by the reconstruction isomorphism. The local-rule corollary applies
that theorem to the original local constructor. The occurrence wrappers
supply the already proved full labels directly. These theorems assume a
real rule on the small source, and construct its whole-graph step.

\lean{Copying/CompleteMatchRewrite}{18}{19}{CompleteMatchCut.replaceWith}
\lean{Copying/CompleteMatchRewrite}{22}{26}{CompleteMatchCut.replacePort}
\lean{Copying/CompleteMatchRewrite}{30}{41}{CompleteMatchCut.rewrite}
\lean{Copying/CompleteMatchRewrite}{45}{54}{CompleteMatchCut.local_rewrite}
\lean{Copying/CompleteMatchRewrite}{57}{61}{Opening.Occurrence.rewrite}
\par\penalty0
\lean{Copying/CompleteMatchRewrite}{63}{67}{Opening.Occurrence.local_rewrite}

\Needspace{8\baselineskip}
\paragraph{Extracting full local occurrences.}
The complete-match occurrence constructor derives the internal geometry
from the incidence laws and full recursive labels. The local-occurrence
theorem inducts over the actual rewrite derivation, composing the occurrence
through every renaming, plug and expansion. The source-token equation is
transported by the same port map at each step.

\lean{Copying/LocalRewriteOccurrence}{13}{31}{MatchEmbedding.occurrence}
\lean{Copying/LocalRewriteOccurrence}{35}{52}{Combined.Rewrite.local_occurrence}

\Needspace{8\baselineskip}
\paragraph{Factoring a supplied rewrite.}
The successor lemma compares the supplied actual step with the newly
realized step using proved successor coherence. \texttt{CutFactorization}
records the actual local rule, full occurrence, exact source token and
pointed successor isomorphism. Its rewrite projection constructs the
canonical successor at the supplied token. The remaining-run equivalence
transports every exact successful length through the successor isomorphism.
The final existence theorem extracts a local occurrence from the supplied
rewrite and assembles these fields.

\lean{Copying/RewriteCutFactorization}{19}{29}{Opening.Occurrence.rewrite_successor}
\lean{Copying/RewriteCutFactorization}{34}{51}{Combined.CutFactorization}
\lean{Copying/RewriteCutFactorization}{54}{61}{Combined.CutFactorization.rewrite}
\lean{Copying/RewriteCutFactorization}{64}{70}{Combined.CutFactorization.terminatesIn_iff}
\lean{Copying/RewriteCutFactorization}{74}{79}{Combined.Rewrite.cut_factorization}
\par\penalty0

\Needspace{8\baselineskip}
\paragraph{Reflection across arbitrary fillings.}
A non-hole certificate has retained contributions on both sides by the
original certificate case split. \texttt{AvoidsShallow} requires this for
exactly the matched edges. The chosen left/right retained corollas yield
the source/target equations at original host coordinates. Incidence coverage
then proves every matched port has a host preimage; \texttt{port} chooses it
and \texttt{host\_port} records the equation. Owners are injective because
two matched edges in one retained corolla would be equal. \texttt{reflect}
uses these coordinates and incidences to construct the common-context
match, and \texttt{across} inserts the opposite retained corollas. Its token
equation follows from injectivity of the original host inclusion.

\lean{Copying/RetainedFootprint}{11}{23}{DeepFillingEdge.retained_of_not_hole}
\lean{Copying/RetainedFootprint}{27}{30}{RetainedFootprint.AvoidsShallow}
\lean{Copying/RetainedFootprint}{32}{33}{RetainedFootprint.left}
\lean{Copying/RetainedFootprint}{35}{36}{RetainedFootprint.right}
\lean{Copying/RetainedFootprint}{40}{45}{RetainedFootprint.source}
\par\penalty0
\lean{Copying/RetainedFootprint}{47}{52}{RetainedFootprint.target}
\lean{Copying/RetainedFootprint}{56}{70}{RetainedFootprint.port_origin}
\lean{Copying/RetainedFootprint}{72}{73}{RetainedFootprint.port}
\lean{Copying/RetainedFootprint}{75}{77}{RetainedFootprint.host_port}
\lean{Copying/RetainedFootprint}{79}{88}{RetainedFootprint.owner_injective}
\par\penalty0
\lean{Copying/RetainedFootprint}{93}{115}{RetainedFootprint.reflect}
\lean{Copying/RetainedFootprint}{119}{132}{RetainedFootprint.across}
\lean{Copying/RetainedFootprint}{135}{139}{RetainedFootprint.across_host}

\Needspace{8\baselineskip}
\paragraph{The original-label matched skeleton.}
A retained contribution has one edge and all its ordered boundary ports,
so it is a labelled corolla; its signature lemmas identify its original
context interface. Full occurrence labels then give the source-label and
signature equations. \texttt{skeleton} retains the matched source carriers
and incidences but assigns the original common-context edge labels.
\texttt{skeletonOccurrence} is a complete occurrence in that same context.
Restrict the original two filling families to its edges; their left copies
collapse to the actual local source, proving \texttt{reconstruction} and
its exact host-token equation.

The original certificate case split restricts each recursive certificate
to the skeleton. Focus freedom is inherited edgewise. No shallow skeleton
edge is a hole, so shallow binding freedom follows immediately; recursive
binding freedom is inherited from the original boxed bodies. The resulting
\texttt{witness} pairs the unchanged actual local source with the opposite
filled skeleton, and its context is proved binding-free. This constructs
a paired footprint; it does not by itself say that the opposite footprint
is the source of a local rule.

\lean{Copying/RetainedFootprintNormalisation}{17}{22}{Search.Expansion.Retained.copy}
\lean{Copying/RetainedFootprintNormalisation}{24}{25}{Search.Expansion.Retained.input}
\lean{Copying/RetainedFootprintNormalisation}{27}{37}{Search.Expansion.Retained.output}
\lean{Copying/RetainedFootprintNormalisation}{41}{45}{RetainedFootprint.sourceLabel}
\lean{Copying/RetainedFootprintNormalisation}{47}{51}{RetainedFootprint.skeleton_input}
\par\penalty0
\lean{Copying/RetainedFootprintNormalisation}{53}{57}{RetainedFootprint.skeleton_output}
\lean{Copying/RetainedFootprintNormalisation}{61}{63}{RetainedFootprint.skeleton}
\lean{Copying/RetainedFootprintNormalisation}{66}{74}{RetainedFootprint.skeletonOccurrence}
\lean{Copying/RetainedFootprintNormalisation}{76}{76}{RetainedFootprint.leftFamily}
\lean{Copying/RetainedFootprintNormalisation}{77}{77}{RetainedFootprint.rightFamily}
\par\penalty0
\lean{Copying/RetainedFootprintNormalisation}{79}{82}{RetainedFootprint.leftCopy}
\lean{Copying/RetainedFootprintNormalisation}{86}{100}{RetainedFootprint.reconstruction}
\lean{Copying/RetainedFootprintNormalisation}{103}{106}{RetainedFootprint.reconstruction_host}
\lean{Copying/RetainedFootprintNormalisation}{109}{115}{RetainedFootprint.skeleton_edges}
\lean{Copying/RetainedFootprintNormalisation}{117}{118}{RetainedFootprint.skeleton_focusFree}
\par\penalty0
\lean{Copying/RetainedFootprintNormalisation}{122}{125}{RetainedFootprint.skeleton_shallowBindingFree}
\lean{Copying/RetainedFootprintNormalisation}{129}{133}{RetainedFootprint.skeleton_bindingFree}
\lean{Copying/RetainedFootprintNormalisation}{137}{145}{RetainedFootprint.witness}
\lean{Copying/RetainedFootprintNormalisation}{147}{149}{RetainedFootprint.witness_bindingFree}

\Needspace{8\baselineskip}
\paragraph{Primitive whole-graph transfer.}
On a primitive footprint, the primitive-matching equation supplies full
label agreement at every edge. Reflecting across the retained contributions
therefore satisfies the complete-match realization theorem. The resulting
opposite-side rewrite starts at the exact paired host token. The local
corollary covers the original primitive local-rule footprints with their
original premises. No shallow-overlap exclusion or successor pairing is
silently inferred: the former is explicit and the latter remains separate.

\lean{Copying/RetainedRewriteTransfer}{14}{28}{MatchEmbedding.primitive_labels}
\lean{Copying/RetainedRewriteTransfer}{32}{40}{RetainedFootprint.primitive_rewrite}
\lean{Copying/RetainedRewriteTransfer}{44}{49}{RetainedFootprint.primitive_local}

\Needspace{8\baselineskip}
\paragraph{The first-hole barrier.}
The shallow projection extracts the outer binding condition from the
recursive definition. Appending one adjacency makes an ancestry path
positive. The head lemma excludes a filled-hole origin by non-entry, and
the adjacency lemma uses the retained predecessor's actual target list.
Induction proves the ancestry-retention lemma: a first hole would violate
the positive binding-path condition. The contraction corollary supplies
its actual source and binding label.

\lean{Copying/BindingFreeOverlap}{11}{15}{BindingFree.shallow}
\lean{Copying/BindingFreeOverlap}{18}{28}{Ancestor.append_adjacent}
\lean{Copying/BindingFreeOverlap}{31}{42}{BindingOverlap.head_retained}
\lean{Copying/BindingFreeOverlap}{46}{56}{BindingOverlap.adjacent_owner}
\lean{Copying/BindingFreeOverlap}{60}{77}{BindingOverlap.ancestor_retained}
\par\penalty0
\lean{Copying/BindingFreeOverlap}{80}{89}{BindingOverlap.contraction_descendant}

\Needspace{8\baselineskip}
\paragraph{Counted copying provenance.}
The trace-to-ancestry lemma forgets only the numerical count. Payload
edge classification is a case split on the normal form, and its ancestry
lemma uses the single lambda-to-box incidence when needed. Each retained
sharing hop reflects its forward label and adjacency; induction preserves
the exact trace length. Source uniqueness identifies the resolution start
with the active contraction, so the full resolution reflects to the common
context. Composing that path with the payload ancestry proves retention
of every payload edge. If its owner were a hole, either it would equal the
contraction owner and have incompatible labels, or the positive ancestry
would violate binding freedom.

\lean{Copying/BindingCopyOverlap}{11}{17}{CopyTrace.ancestor}
\lean{Copying/BindingCopyOverlap}{19}{35}{CopyRule.NormalForm.edge_cases}
\lean{Copying/BindingCopyOverlap}{38}{51}{CopyRule.NormalForm.head_ancestor}
\lean{Copying/BindingCopyOverlap}{54}{67}{BindingOverlap.copying_hop_owner}
\lean{Copying/BindingCopyOverlap}{70}{84}{BindingOverlap.copying_trace_owner}
\par\penalty0
\lean{Copying/BindingCopyOverlap}{87}{105}{BindingOverlap.copying_resolution_owner}
\lean{Copying/BindingCopyOverlap}{108}{120}{BindingOverlap.copying_payload_retained}
\lean{Copying/BindingCopyOverlap}{123}{138}{BindingOverlap.copying_payload_not_hole}

\Needspace{8\baselineskip}
\paragraph{Admissible paired fillings.}
The specimen constructor keeps the shared context and certificate families,
transporting endpoint hole freedom through their isomorphisms. Its left and
right class equations are pointed recursive isomorphisms. The quasi theorem
uses those equations to construct the original binding-free specimen under
the clause's rootedness premises; the initial-token corollary specializes
to input zero. A literal pair uses the primitive-hole binding proof.
Serial assembly has its literal serial context; its binding proof requires
the stated cross-path exclusion. Parallel and tensor assembly use independent
components. Boundary reindexing preserves the context after eliminating
its equality proof. Lambda assembly retains the boxed body and has no
shallow holes. Copying preserves this shielding on the retained branch,
uses parallel binding freedom for the two copies, and appends the closed
distributors. The deep-copy corollary supplies the lambda shielding proof.

\lean{Copying/BindingFillingWitness}{13}{27}{FillingWitness.bindingSpecimen}
\lean{Copying/BindingFillingWitness}{29}{33}{FillingWitness.bindingSpecimen_left}
\lean{Copying/BindingFillingWitness}{35}{39}{FillingWitness.bindingSpecimen_right}
\lean{Copying/BindingFillingWitness}{43}{57}{FillingWitness.binding_quasi}
\lean{Copying/BindingFillingWitness}{59}{65}{FillingWitness.binding_initial_quasi}
\par\penalty0
\lean{Copying/BindingFillingWitness}{69}{71}{FillingWitness.of_rel_binding}
\lean{Copying/BindingFillingWitness}{74}{81}{FillingWitness.serial_context}
\lean{Copying/BindingFillingWitness}{83}{90}{FillingWitness.serial_binding}
\lean{Copying/BindingFillingWitness}{92}{97}{FillingWitness.parallel_binding}
\lean{Copying/BindingFillingWitness}{100}{108}{FillingWitness.tensor_binding}
\par\penalty0
\lean{Copying/BindingFillingWitness}{110}{123}{FillingWitness.reboundary_binding}
\lean{Copying/BindingFillingWitness}{125}{127}{FillingWitness.interfaces_binding}
\lean{Copying/BindingFillingWitness}{131}{137}{FillingWitness.lambdaPayload_binding}
\lean{Copying/BindingFillingWitness}{139}{145}{FillingWitness.lambdaPayload_no_shallow_holes}
\lean{Copying/BindingFillingWitness}{150}{171}{FillingWitness.copyingAfter_binding}
\par\penalty0
\lean{Copying/BindingFillingWitness}{175}{181}{FillingWitness.deepCopy_binding}

\Needspace{8\baselineskip}
\paragraph{Admissible host insertion.}
The equality-cast lemma records that casting endpoint graph indices leaves
the stored context unchanged. The plug-context isomorphism follows the
selected component through the dependent family and uses the original
single-expansion isomorphism. Sourced-hole closure then proves binding
freedom of the plugged witness. \texttt{plugHost} retains an original host
port. The two single-expansion position equations follow both flattenings
and their endpoint maps. The final quasi theorem uses expansion peeling
to obtain the actual plugged endpoint graphs and the original typed resume
coordinate, then constructs the binding-free specimen.

\lean{Copying/BindingFillingPlug}{13}{20}{FillingWitness.cast_context}
\lean{Copying/BindingFillingPlug}{22}{40}{FillingWitness.plug_context_iso}
\lean{Copying/BindingFillingPlug}{44}{50}{FillingWitness.plug_binding}
\lean{Copying/BindingFillingPlug}{54}{61}{FillingWitness.plugHost}
\lean{Copying/BindingFillingPlug}{63}{68}{FillingWitness.singleExpansion_leftPosition}
\par\penalty0
\lean{Copying/BindingFillingPlug}{70}{75}{FillingWitness.singleExpansion_rightPosition}
\lean{Copying/BindingFillingPlug}{79}{117}{FillingWitness.plug_binding_quasi}

\Needspace{8\baselineskip}
\paragraph{Box-context bookkeeping.}
Boundary reindexing leaves the common graph context unchanged up to the
same boundary cast; the proof eliminates the equality witnesses. Rebuilt
boxed copies consequently have one original box corolla whose body is the
actual common body context. This explicit label-copy certificate is the
recursive invariant required by padding.

\lean{Copying/BindingFillingBoxes}{10}{23}{FillingWitness.reboundary_context}
\lean{Copying/BindingFillingBoxes}{27}{34}{FillingWitness.boxedCopies_contextCopy}

\Needspace{8\baselineskip}
\paragraph{Padding with certified context geometry.}
The related-corolla collapse theorem uses the existing carrier bijections
and ordered incidences, with recursive label equivalence in place of label
equality. \texttt{CorollaContext} records that a component's common context
is one such corolla; boundary casting preserves it. Retained copies and
literal pairs satisfy this invariant, with the literal pair retaining its
original hole name and type. Expanded padding collapses these component
contexts to the padded original context. The hole case uses template
padding closure and the signature casts.

The recursive existence proof handles retained edges, actual filled holes
and original boxes separately. In the box case it extracts the body's
explicit expansion data and applies the induction hypothesis. It returns
both a padded filling witness and a recursive isomorphism from its actual
context to the padded original context. \texttt{raiseWithContext} chooses
this strengthened pair of data. \texttt{raiseBinding} transports binding
freedom along that certified isomorphism and the proved padding law. It
makes no assertion about the unrelated witness selected by an earlier,
unqualified choice.

\lean{Copying/BindingFillingPadding}{14}{34}{FillingCollapse.related}
\lean{Copying/BindingFillingPadding}{38}{42}{FillingWitness.CorollaContext}
\lean{Copying/BindingFillingPadding}{44}{49}{FillingWitness.CorollaContext.reboundary}
\lean{Copying/BindingFillingPadding}{51}{54}{FillingWitness.of_copies_context}
\lean{Copying/BindingFillingPadding}{56}{58}{FillingWitness.of_rel_context}
\par\penalty0
\lean{Copying/BindingFillingPadding}{62}{71}{FillingWitness.paddedExpansion_context}
\lean{Copying/BindingFillingPadding}{74}{88}{FillingWitness.paddedHole_context}
\lean{Copying/BindingFillingPadding}{93}{147}{FillingWitness.raisedContext_exists}
\lean{Copying/BindingFillingPadding}{150}{154}{FillingWitness.raiseWithContext}
\lean{Copying/BindingFillingPadding}{158}{163}{FillingWitness.raiseBinding}
\par\penalty0

\Needspace{8\baselineskip}
\paragraph{The directed beta splice.}
The two splice slots have no self-loop, by their explicit input/output
ports. An argument-to-body adjacency would identify an external output
with the external body input, so is impossible. Induction makes the slot
index monotone along a path; equal endpoint slots force the whole path to
stay in one component. A forbidden binding path ending at a hole must
therefore stay in the body, since the argument has no holes. Component
binding freedom excludes it. The final theorem combines this shallow
argument with focus freedom and recursive binding freedom of every edge.

\lean{Copying/BindingSplice}{12}{25}{BindingFree.splice_nonSelfLoop}
\lean{Copying/BindingSplice}{28}{38}{BindingFree.splice_no_backward}
\lean{Copying/BindingSplice}{41}{50}{BindingFree.splice_ancestor_monotone}
\lean{Copying/BindingSplice}{53}{73}{BindingFree.splice_ancestor_same}
\lean{Copying/BindingSplice}{76}{107}{BindingFree.shallow_splice_suffix}
\par\penalty0
\lean{Copying/BindingSplice}{111}{128}{BindingFree.splice_suffix}

\Needspace{8\baselineskip}
\paragraph{The beta residual.}
The splice-context equation exposes the actual shared splice assembled by
the witness. The splice-binding theorem applies the directed-splice lemma.
A stable paired argument is recursively isomorphic to its other endpoint,
so \texttt{stableArgument} represents both endpoints using the actual stable
graph as common context; stability supplies its hole freedom. The
\texttt{betaResidual} constructor pads the body with certified context
geometry, casts its ordered bound/free interface, and splices that stable
argument. Its subtype result includes binding freedom of the constructed
common context for the two original beta successors.

\lean{Copying/BindingFillingBeta}{15}{22}{FillingWitness.splice_context}
\lean{Copying/BindingFillingBeta}{26}{36}{FillingWitness.splice_binding}
\lean{Copying/BindingFillingBeta}{41}{44}{OperationalTemplates.DeepTransfer.stableArgument}
\lean{Copying/BindingFillingBeta}{46}{48}{OperationalTemplates.DeepTransfer.stableArgument_closed}
\lean{Copying/BindingFillingBeta}{53}{64}{OperationalTemplates.DeepTransfer.betaResidual}
\par\penalty0

\Needspace{8\baselineskip}
\paragraph{Operational invariants.}
The pair-raising theorem transports the original literal rules and endpoint
isomorphisms; the padding-closed theorem exposes precisely this property.
The beta reduct cannot have only stable labels because its retained body
still has a term output. The literal nonstability theorem combines that
fact with the active source heads and the reference reduct's store; the
pair theorem transports it through endpoint isomorphisms. Stable-label
preservation follows because neither operational side is flat stable.
A stable edge certificate can thus only be retained; the filled and boxed
cases contradict stability. Collapsing all those retained corollas proves
the expansion isomorphism and hence the endpoint witness isomorphism.

\lean{Copying/OperationalTemplateLaws}{10}{22}{OperationalTemplates.Pair.raise}
\lean{Copying/OperationalTemplateLaws}{24}{24}{OperationalTemplates.template_paddingClosed}
\lean{Copying/OperationalTemplateLaws}{28}{34}{OperationalTemplates.beta_after_not_stable_labels}
\lean{Copying/OperationalTemplateLaws}{38}{55}{OperationalTemplates.Literal.not_stable_labels}
\lean{Copying/OperationalTemplateLaws}{57}{62}{OperationalTemplates.Pair.not_stable_labels}
\par\penalty0
\lean{Copying/OperationalTemplateLaws}{64}{67}{OperationalTemplates.template_preservesStableLabels}
\lean{Copying/OperationalTemplateLaws}{71}{82}{OperationalTemplates.stable_edge_retained}
\lean{Copying/OperationalTemplateLaws}{84}{90}{OperationalTemplates.stable_expansion_iso}
\lean{Copying/OperationalTemplateLaws}{94}{97}{OperationalTemplates.stable_witness_iso}

\Needspace{8\baselineskip}
\paragraph{Whole-host deep diagrams.}
The generic framed theorem plugs both already proved original rewrites
into the same sourced-hole binding-free closing host. Their target hole
freedom follows from actual rewrite preservation, and the proved residual
witness constructs the original binding-free quasi specimen. The input
index is the unchanged numeric index under an ordered interface cast;
its typing and rewrite-cast lemmas follow by eliminating the interface
equalities. The copying theorem supplies the original local copying rules
and the constructed duplicated residual. The beta theorem supplies the
original application rules, derives the opposite stable argument from
its actual isomorphism, and uses the certified beta residual. Thus both
whole-host diagrams include actual transitions and admissible pointed
residuals, without assuming template robustness.

\lean{Copying/OperationalDeepTransfer}{15}{40}{FillingWitness.binding_framed_rewrite}
\lean{Copying/OperationalDeepTransfer}{42}{43}{OperationalTemplates.DeepTransfer.inputIndex}
\lean{Copying/OperationalDeepTransfer}{45}{50}{OperationalTemplates.DeepTransfer.rewrite_input_cast}
\lean{Copying/OperationalDeepTransfer}{52}{55}{OperationalTemplates.DeepTransfer.inputIndex_term}
\lean{Copying/OperationalDeepTransfer}{60}{90}{OperationalTemplates.DeepTransfer.copy_framed_diamond}
\par\penalty0
\lean{Copying/OperationalDeepTransfer}{96}{128}{OperationalTemplates.DeepTransfer.beta_framed_diamond}

\Needspace{8\baselineskip}
\paragraph{Operational regression certificates.}
\texttt{unusedBody} tensors the original graph with a term weakening;
\texttt{unusedWitness} tensors its original certificate with the reflexive
weakening certificate. Independent tensor binding freedom proves the
unused-binder lemma. The reference and beta witnesses place the actual
operational pairs in primitive holes and apply this construction. Their
binding lemmas give admissible deep-copy residuals. The reference step
lemma separately checks the original copy transition on each side.
The nested reference witness adds a second lambda and unused output;
box and tensor closure prove its binding freedom, which the same copy
residual theorem preserves in both branches.

\lean{Copying/OperationalDeepExamples}{11}{13}{OperationalDeep.Examples.unusedBody}
\lean{Copying/OperationalDeepExamples}{15}{18}{OperationalDeep.Examples.unusedWitness}
\lean{Copying/OperationalDeepExamples}{20}{25}{OperationalDeep.Examples.unused_binding}
\lean{Copying/OperationalDeepExamples}{27}{28}{OperationalDeep.Examples.referenceWitness}
\lean{Copying/OperationalDeepExamples}{30}{31}{OperationalDeep.Examples.reference_binding}
\par\penalty0
\lean{Copying/OperationalDeepExamples}{35}{37}{OperationalDeep.Examples.reference_copy_binding}
\lean{Copying/OperationalDeepExamples}{40}{55}{OperationalDeep.Examples.reference_copy_steps}
\lean{Copying/OperationalDeepExamples}{57}{58}{OperationalDeep.Examples.betaWitness}
\lean{Copying/OperationalDeepExamples}{60}{61}{OperationalDeep.Examples.beta_binding}
\lean{Copying/OperationalDeepExamples}{64}{66}{OperationalDeep.Examples.beta_copy_binding}
\par\penalty0
\lean{Copying/OperationalDeepExamples}{69}{70}{OperationalDeep.Examples.nestedReferenceWitness}
\lean{Copying/OperationalDeepExamples}{72}{73}{OperationalDeep.Examples.nestedReference_binding}
\lean{Copying/OperationalDeepExamples}{77}{79}{OperationalDeep.Examples.nestedReference_copy_binding}

\Needspace{8\baselineskip}
\paragraph{One common frame for both fillings.}
Cut the fully labelled skeleton occurrence out of the original common
context. The grouped family puts the selected skeleton expansion at the
new hole and leaves every unmatched original filling independent. An edge
partition separates matched owners from their complement. The vertex and
edge split maps merely regroup the resulting sum and dependent-pair tags;
their inverses undo that regrouping. Combining them with the cut carrier
bijections gives the whole expansion bijections.

The host identity follows by port cases. The local-host identity uses the
original occurrence's internal map and the cut boundary equations. The
retained source/target-port identities follow from the projected original
lists; they give the retained insertion identity. Applying the skeleton's
ordered incidence equations at a finite index gives its source/target-port
identities and hence the selected insertion identity. The final recursive
isomorphism checks labels and each ordered incidence on the two edge cases.
It reconstructs the entire original expansion for any filling family, so
the same frame serves both sides without restricting unrelated multi-edge
fillings. Binding freedom of this auxiliary cut is not asserted.

\lean{Copying/RetainedFootprintFrame}{17}{22}{RetainedFootprint.frame}
\lean{Copying/RetainedFootprintFrame}{24}{24}{RetainedFootprint.localFamily}
\lean{Copying/RetainedFootprintFrame}{26}{32}{RetainedFootprint.grouped}
\lean{Copying/RetainedFootprintFrame}{34}{36}{RetainedFootprint.frameEdgePartition}
\lean{Copying/RetainedFootprintFrame}{39}{59}{RetainedFootprint.frameVertexSplit}
\par\penalty0
\lean{Copying/RetainedFootprintFrame}{61}{65}{RetainedFootprint.frameVertices}
\lean{Copying/RetainedFootprintFrame}{67}{77}{RetainedFootprint.frameEdgeSplit}
\lean{Copying/RetainedFootprintFrame}{79}{82}{RetainedFootprint.frameEdges}
\lean{Copying/RetainedFootprintFrame}{84}{88}{RetainedFootprint.frame_host}
\lean{Copying/RetainedFootprintFrame}{90}{111}{RetainedFootprint.frame_local_host}
\par\penalty0
\lean{Copying/RetainedFootprintFrame}{113}{117}{RetainedFootprint.frame_sourcePort}
\lean{Copying/RetainedFootprintFrame}{119}{123}{RetainedFootprint.frame_targetPort}
\lean{Copying/RetainedFootprintFrame}{125}{137}{RetainedFootprint.frame_other}
\lean{Copying/RetainedFootprintFrame}{139}{147}{RetainedFootprint.skeleton_sourcePort}
\lean{Copying/RetainedFootprintFrame}{149}{157}{RetainedFootprint.skeleton_targetPort}
\par\penalty0
\lean{Copying/RetainedFootprintFrame}{159}{176}{RetainedFootprint.frame_local_insert}
\lean{Copying/RetainedFootprintFrame}{180}{221}{RetainedFootprint.frameIso}

\Needspace{8\baselineskip}
\paragraph{Quantitative deep robustness clauses.}
The forward generic wrapper consumes the two actual steps and their
constructed binding-free quasi specimen. Proved class determinism
identifies any independently supplied left successor with the constructed
one. The cost witnesses are zero additional steps on that side and one
step on the opposite side, so reflexivity proves the required comparison
of one with one for every cost order. The converse exchanges the two
specimen projections and uses the same argument; it does not reverse a
compute rule. The framed source definition records the actual host source
port in rewrite mode. The four concrete corollaries instantiate these
wrappers with the proved copying or beta host diagram, in each direction.
They quantify over arbitrary supplied successors once the framed body and
context data are constructed. Deriving those data from every arbitrary
original specimen remains outside these case theorems.

\lean{Copying/OperationalDeepRobustness}{12}{24}{FillingWitness.binding_framed_robust_case}
\lean{Copying/OperationalDeepRobustness}{29}{44}{FillingWitness.binding_framed_robust_case_opposite}
\lean{Copying/OperationalDeepRobustness}{46}{63}{OperationalTemplates.DeepTransfer.framedSource}
\lean{Copying/OperationalDeepRobustness}{67}{78}{OperationalTemplates.DeepTransfer.copy_robust_case}
\lean{Copying/OperationalDeepRobustness}{81}{108}{OperationalTemplates.DeepTransfer.copy_robust_case_opposite}
\par\penalty0
\lean{Copying/OperationalDeepRobustness}{112}{121}{OperationalTemplates.DeepTransfer.beta_robust_case}
\lean{Copying/OperationalDeepRobustness}{125}{134}{OperationalTemplates.DeepTransfer.beta_robust_case_opposite}

\Needspace{8\baselineskip}
\paragraph{Grafting the local residual into the whole graph.}
The edge-witness constructor turns each retained, filled or boxed
certificate into an actual paired component witness. The successor family
changes only the selected grouped component; its graph and port definitions
use the actual common-frame expansion and insertion. The frame rewrite
lifts the supplied local step by the original expansion constructor, then
uses the grouping isomorphism. The inverse reconstruction equation
identifies the left local token, proving the left whole-graph rewrite at
its original occurrence coordinate.

The residual witness expands the same frame with the proved local residual
at the selected component and witnesses for the original unmatched endpoint pairs everywhere
else. The paired-rewrite theorem assembles the two actual whole-graph
steps and this genuine recursive filling certificate. The precise right
token theorem uses its explicit attachment equality; the independent left
successor is identified with the constructed successor by proved
coherence. This establishes the whole-graph and residual-certificate part
of transfer. Binding freedom of the grafted shared context is a separate
path obligation; original hole names and pointed residual-coordinate
correspondence also remain separate. They are not supplied by a new locality
assumption.

\lean{Copying/RetainedFrameTransfer}{15}{33}{FillingWitness.of_edge}
\lean{Copying/RetainedFrameTransfer}{37}{45}{RetainedFootprint.afterFamily}
\lean{Copying/RetainedFrameTransfer}{47}{50}{RetainedFootprint.afterGraph}
\lean{Copying/RetainedFrameTransfer}{52}{57}{RetainedFootprint.afterPort}
\lean{Copying/RetainedFrameTransfer}{61}{79}{RetainedFootprint.frame_rewrite}
\par\penalty0
\lean{Copying/RetainedFrameTransfer}{81}{84}{RetainedFootprint.reconstruction_inverse}
\lean{Copying/RetainedFrameTransfer}{87}{100}{RetainedFootprint.left_rewrite}
\lean{Copying/RetainedFrameTransfer}{104}{115}{RetainedFootprint.right_rewrite_at}
\lean{Copying/RetainedFrameTransfer}{119}{124}{RetainedFootprint.left_successor}
\lean{Copying/RetainedFrameTransfer}{128}{133}{RetainedFootprint.residual}
\par\penalty0
\lean{Copying/RetainedFrameTransfer}{138}{151}{RetainedFootprint.paired_rewrites}

